\chapter{Porting Revizor to AArch64}
\label{chap:port}

This chapter describes the first contribution: bringing model-based relational testing to
AArch64. \revizor{}'s pipeline is ISA-independent, so the structure is kept and the
ISA-specific backends are rebuilt for AArch64, with new support for the Arm security
extensions and for fuzzing physically remote cores. \autoref{sec:port-overview} gives the
overall shape, \autoref{sec:port-components} the native-execution architecture,
\autoref{sec:port-isa} how the instruction set is derived from Arm's specification,
\autoref{sec:port-ce} the contract executor, \autoref{sec:port-features} the modeling of
Pointer Authentication and Memory Tagging, \autoref{sec:port-driving} the executor controls,
\autoref{sec:port-input} the input format, and \autoref{sec:port-remote} remote fuzzing.

\section{Overview}
\label{sec:port-overview}

\revizor{}~\cite{revizor} is organized as a pipeline (a generator, a model, an executor,
and an analyzer) in which the analyzer compares the model's contract trace against the
executor's hardware trace (\autoref{sec:bg-mrt}). This structure is independent of the
instruction set, so we keep it and rebuild the three ISA-specific backends for AArch64, the
64-bit execution state introduced in Armv8-A. We additionally model the Pointer
Authentication and Memory Tagging extensions, which x86 \revizor{} has no analogue for and
which are the subject of the precondition primitives of \autoref{chap:precond}. The port is
evaluated on Armv9.2-A cores in \autoref{chap:eval}.

\section{Native execution on the real core}
\label{sec:port-components}

Mainline \revizor{} runs as a single process, with the model on an in-process software
emulator of the contract. Faithfully emulating a modern AArch64 core in software is both
laborious and a source of divergence from real silicon. We therefore run the model as a
separate process that executes the test case \emph{natively} on a real core. The fuzzer then
coordinates two external components, each over its own channel (\autoref{fig:port-components}).
The \emph{contract executor}
(CE) produces the contract trace. A Linux \emph{kernel module}, reached through a
\reg{/dev/executor} character device, runs the test case on the CPU under test and measures
its microarchitectural footprint. The fuzzer generates programs and inputs, drives the CE
and the kernel module, and compares the traces they return. Because the kernel module is the
only component that must sit on the silicon under test, it can run on a different machine
from the rest of the fuzzer, which is what makes remote fuzzing (\autoref{sec:port-remote})
possible.

\begin{figure}[t]
\centering
\fitwidth{\begin{tikzpicture}[
  font=\footnotesize,
  box/.style={draw=black!70, fill=black!4, rounded corners=2.5pt, thick,
              align=center, inner sep=3pt},
  ch/.style={font=\scriptsize, text=black!70, align=center, inner sep=1.5pt},
  ar/.style={-{Stealth[length=1.8mm,width=1.6mm]}, thick, draw=black!75},
]
  \node[box, minimum width=116mm, minimum height=8mm] (fz) at (0,0)
    {\textbf{Fuzzer}~\textnormal{(Python)}\\[-1pt]%
     {\scriptsize generate \,$\to$\, model \,$\to$\, measure \,$\to$\, compare}};
  \node[box, minimum width=42mm, minimum height=11mm] (ce) at (-3.1,-2.6)
    {Contract\\executor (CE)\\[-1pt]{\scriptsize model $\to$ ctrace}};
  \node[box, minimum width=44mm, minimum height=11mm] (km) at (3.1,-2.6)
    {Kernel module\\{\scriptsize \texttt{/dev/executor}}\\[-1pt]{\scriptsize run on CUT $\to$ htrace}};
  \node[box, dashed, minimum width=42mm, minimum height=6mm] (bpu) at (-3.1,-4.6)
    {{\scriptsize CPython: execution clause}};
  \node[box, minimum width=44mm, minimum height=6mm] (cpu) at (3.1,-4.6)
    {CPU {\scriptsize (CUT)}};
  \draw[dashed, rounded corners=3pt, draw=black!55] (0.75,-1.87) rectangle (5.45,-5.08);
  \node[font=\scriptsize\itshape, text=black!55] at (3.1,-5.4) {target device (local or remote)};
  \draw[ar] (ce.north) -- (ce.north |- fz.south) node[ch,midway,left=1pt]{pipe};
  \draw[ar] (km.north) -- (km.north |- fz.south)
        node[ch,midway,right=1pt]{\texttt{ioctl}, \texttt{sysfs}};
  \draw[ar] (ce) -- (bpu);
  \draw[ar] (km) -- (cpu);
\end{tikzpicture}
}
\caption{Component map. The fuzzer coordinates the contract executor and the kernel module,
each over its own channel: the CE returns the contract trace, and the kernel module runs the
test case on the CPU under test and returns the hardware trace. Only the kernel module must
sit on the silicon under test, so it may run on a remote target device.}
\label{fig:port-components}
\end{figure}

\section{Deriving the instruction set from the Arm specification}
\label{sec:port-isa}

The generator needs a precise model of every instruction it may emit: its encoding,
operands, the registers and memory it reads and writes, and the legal domains of its
immediates. Mainline \revizor{} reads this table from a prebuilt file. We instead derive it
directly from Arm's machine-readable specification, taking each instruction's encoding,
operand roles, register and memory access roles, and immediate domains from its Architecture
Specification Language (ASL) definition. Encodings we cannot fully recover from the
specification are reported and excluded, so the generator emits only
instructions that have a complete, validated model. Pinning the specification release keeps
the instruction set reproducible across fuzzing campaigns.

\section{The contract executor: native single-stepping}
\label{sec:port-ce}

Rather than emulate AArch64, the CE runs the test case on the real core and observes it one
instruction at a time (\autoref{fig:port-trampoline}). It is a user-space process that loads the
program into an executable
mapping and single-steps it. The single-stepping exploits the fixed-width A64 encoding: every
instruction is 32 bits wide, and so is a \reg{BL} (branch-with-link). Before the test case runs, the CE replaces each and every
instruction word in place with a \reg{BL} to a trampoline. A \reg{BL} does not modify the
condition flags, so the \reg{NZCV} state carries across every interposed branch into the next
test instruction. To run a single instruction the CE restores that one word to its original
encoding, leaves every other word as its trampoline branch, and transfers control to it. The
core executes that one instruction natively and then continues along the program's own control
flow to the next instruction it reaches, which need not be the one immediately after it, since a
branch or a loop may redirect control. That next word is still a \reg{BL}, so execution falls
into the trampoline, where a hook records the architectural state. The CE then restores the next
instruction and reinstates the trampoline at the one just executed, so that exactly one original
instruction is ever in place, and repeats until the test case completes.

Running on the CPU under test avoids the emulation inconsistencies we observed in common
emulators such as Unicorn~\cite{unicorn}. It also strikes a middle ground between full
emulation and purely local execution. A complete emulator can model a feature the host lacks,
but risks fidelity gaps. Purely native execution is faithful, but needs the feature present on
the local machine. The CE takes the middle path: it runs base instructions natively and
models, on request, the extensions in software. For example, we fuzzed MTE on a Google
Pixel~8, which implements the extension, while generating and modeling on a host that does
not.

\begin{figure}[t]
\centering
\fitwidth{\begin{tikzpicture}[
  font=\footnotesize,
  stp/.style={draw=black!70, fill=black!4, rounded corners=2.5pt, thick,
              align=center, text width=44mm, minimum height=7mm, inner sep=3pt},
  ar/.style={-{Stealth[length=1.8mm,width=1.6mm]}, thick, draw=black!75},
  lbl/.style={font=\scriptsize, text=black!65, inner sep=1.5pt},
]
  \node[stp] (rs) at (0,0)    {restore instruction $i$ at the PC};
  \node[stp] (ex) at (0,-1.55){execute instruction $i$ natively on the CUT};
  \node[stp] (tr) at (0,-3.1) {\texttt{BL} trampoline: save register context};
  \node[stp] (hk) at (0,-4.95){hook: record architectural state,\\ drive speculation\\ {\scriptsize (checkpoint $\cdot$ redirect $\cdot$ roll back)}};
  \draw[ar] (rs) -- (ex);
  \draw[ar] (ex) -- (tr) node[lbl,midway,right=1pt]{next executed word is \texttt{BL hook}};
  \draw[ar] (tr) -- (hk);
  \draw[ar] (hk.west) -- ++(-0.75,0) |- (rs.west);
  \node[lbl, rotate=90, fill=white, inner sep=1pt] at (-2.6,-2.475) {next instruction};
  \path let \p1 = (current bounding box.north east),
            \p2 = (current bounding box.south west),
            \n1 = {max(\x1, -\x2)} in
        coordinate (bbne) at (\n1, \y1) coordinate (bbsw) at (-\n1, \y2);
  \pgfresetboundingbox
  \useasboundingbox (bbsw) rectangle (bbne);
\end{tikzpicture}
}
\caption{The CE single-steps the test case. Every instruction word is replaced by a \reg{BL} to
the trampoline, and only the next instruction to execute is restored to its original encoding.
After that instruction runs natively, control reaches the next fetched word, which is again a
\reg{BL}: for a non-branch (a) it is the following instruction, and for a taken branch (b) it is
the branch target. The trampoline (\texttt{base\_hook}) saves the register context and calls the
C hook (\texttt{base\_hook\_c}), which snapshots the architectural state and memory, models the
requested speculation, emulates \reg{PAC} and \reg{MTE}, and re-arms the slots so that only the
next instruction is left native.}
\label{fig:port-trampoline}
\end{figure}

\paragraph{The execution clause and the speculation stack.}
A contract's execution clause specifies what speculation the CE must model, and the CE
realizes it with a \emph{speculation stack} (\autoref{fig:port-specstack}). At a speculation
point (a mispredicted branch, a bypassed store), the CE pushes a checkpoint: a snapshot of the registers, the simulated
memory, the MTE tag memory, and the architectural call stack. It then redirects execution
onto the speculative path. When that path ends or exhausts its per-window instruction budget,
the CE pops the checkpoint, restores the snapshot, and resumes at the recorded address.
The speculation windows the CE simulates may nest to a configurable depth, and the checkpoints
form a last-in-first-out stack, so an inner window unwinds before the one enclosing it. The execution clause decides where a window opens,
and its logic runs in the CE's embedded Python interpreter, so a new execution clause plugs in
without touching the C engine.

\begin{figure}[hb]
\centering
\fitwidth{\begin{tikzpicture}[
  font=\footnotesize,
  w/.style={draw=black!70, fill=black!4, rounded corners=2.5pt, thick,
            align=center, minimum width=15mm, minimum height=6.5mm, inner sep=2pt},
  spec/.style={w, dashed, fill=white},
  e/.style={-{Stealth[length=1.7mm,width=1.4mm]}, thick, draw=black!75},
  lbl/.style={font=\scriptsize, text=black!65, inner sep=1.5pt},
  axis/.style={-{Stealth[length=1.9mm]}, draw=black!45},
]
  \draw[axis] (-1.5,0.35) -- (-1.5,-2.85);
  \node[font=\scriptsize, text=black!55, rotate=90, anchor=south] at (-1.78,-1.25) {nesting depth};
  \node[lbl, anchor=east] at (-1.22, 0)     {0};
  \node[lbl, anchor=east] at (-1.22,-1.25)  {1};
  \node[lbl, anchor=east] at (-1.22,-2.5)   {2};
  \node[w]    (w0) at (0,0)        {architectural};
  \node[w]    (a0) at (2.3,0)      {continue};
  \node[spec] (w1) at (1.0,-1.25)  {window};
  \node[spec] (a1) at (3.3,-1.25)  {continue};
  \node[spec] (w2) at (2.0,-2.5)   {window};
  \draw[e] (w0) -- (a0);
  \draw[e] (w0) -- (w1) node[lbl,midway,left=0pt]{spec.};
  \draw[e] (w1) -- (a1);
  \draw[e] (w1) -- (w2) node[lbl,midway,left=0pt]{spec.};
  \path let \p1 = (current bounding box.north east),
            \p2 = (current bounding box.south west),
            \n1 = {max(\x1-1.65cm, 1.65cm-\x2)} in
        coordinate (bbne) at ({1.65cm+\n1}, \y1)
        coordinate (bbsw) at ({1.65cm-\n1}, \y2);
  \pgfresetboundingbox
  \useasboundingbox (bbsw) rectangle (bbne);
\end{tikzpicture}
}
\caption{Execution forks at each speculation point into a speculative window (dashed) and the
continuation. Windows nest to deeper levels and unwind innermost first, so the checkpoints form
a last-in-first-out stack.}
\label{fig:port-specstack}
\end{figure}

\section{Modeling the Arm security extensions}
\label{sec:port-features}

The precondition primitives of \autoref{chap:precond} need the model to reproduce the operand
checks of Pointer Authentication and Memory Tagging faithfully, even on a host that does not
implement them. Pointer Authentication is modeled in software, so the CE can compute and verify a PAC the way
the target core would. The model signs and authenticates under the keys supplied with each
input, and applies the QARMA variant and PAC field layout it is configured to match. Memory Tagging is modeled over a simulated tag memory that the speculation-stack
checkpoints snapshot and restore alongside the data memory, so a simulation of transiently
executed tagged access sees, and can violate, the correct allocation tag. Both extensions are supplied with
per-input state (custom PAC keys and an initial tag memory) through the input file described
next, which lets a single test case be exercised under many keys and tag layouts.

\section{Driving the executor}
\label{sec:port-driving}

Control of the kernel module is a small \reg{ioctl} state machine (\autoref{fig:port-execfsm}),
built around a checkout model. The host \emph{checks out} the test region to load the body, and
\emph{allocates} each input, where every allocation returns an input identifier (\emph{iid}) that
names its state slot. Checking out an \emph{iid} makes that slot the active target, so the next
write fills it, and many inputs can be prepared before a run. A \emph{trace} then runs the loaded
test over all inputs, and a \emph{get-measurement} reads one input's result back by its \emph{iid}.
A few further commands round out the same model (unloading, freeing, clearing, and querying), and a
separate group exposes the PAC primitives (sign, authenticate, strip) that the contract model and
the validation path use. A run is possible only once both a test and its inputs are present, and
the traced state reproduces across re-runs. Each measurement reports, per input, the L1 data-cache
sets the execution touched (the hardware trace of \autoref{sec:bg-mrt}), read through a cache side
channel.

A second group of controls, exposed through \reg{sysfs}, configures how a run measures rather than
what it runs. The knobs that matter in practice are the measurement method (Prime+Probe,
Flush+Reload, or Prime+Reload), the pinned core, the warmup count, and a few toggles over
the cache and branch-history state each run starts from, together with the speculative-store-bypass
setting.

\begin{figure}[t]
\centering
\fitwidth{\begin{tikzpicture}[
  font=\footnotesize,
  st/.style={draw=black!70, fill=black!4, rounded corners=2.5pt, thick,
             align=center, minimum width=24mm, minimum height=9mm, inner sep=2pt},
  fin/.style={st, fill=black!12},
  ed/.style={-{Stealth[length=1.6mm,width=1.4mm]}, semithick, draw=black!75},
  lbl/.style={font=\scriptsize, text=black!70, inner sep=2pt, fill=white},
]
  \node[st]  (cfg) at (0,0)      {\textsc{config}\\{\scriptsize no test, no inputs}};
  \node[st]  (lt)  at (4.0,1.7)  {\textsc{loaded test}\\{\scriptsize test only}};
  \node[st]  (li)  at (4.0,-1.7) {\textsc{loaded inputs}\\{\scriptsize inputs only}};
  \node[st]  (rdy) at (8.0,0)    {\textsc{ready}\\{\scriptsize test + inputs}};
  \node[fin] (tr)  at (12.8,0)   {\textsc{traced}\\{\scriptsize measured}};
  \draw[ed] (cfg) to[bend left=12]  node[lbl,above left=-3pt]{write test}   (lt);
  \draw[ed] (lt)  to[bend left=12]  node[lbl,below right=-3pt]{unload test} (cfg);
  \draw[ed] (cfg) to[bend right=12] node[lbl,below left=-3pt]{write input}  (li);
  \draw[ed] (li)  to[bend right=12] node[lbl,above right=-3pt]{clear inputs} (cfg);
  \draw[ed] (lt)  to[bend left=12]  node[lbl,above right=-3pt]{write input} (rdy);
  \draw[ed] (rdy) to[bend left=12]  node[lbl,below left=-3pt]{clear inputs} (lt);
  \draw[ed] (li)  to[bend right=12] node[lbl,below right=-3pt]{write test}  (rdy);
  \draw[ed] (rdy) to[bend right=12] node[lbl,above left=-3pt]{unload test}  (li);
  \draw[ed] (rdy) -- node[lbl,above]{\textsc{trace}} (tr);
  \draw[ed] (tr) to[out=-60,in=-120,looseness=6] node[lbl,below]{re-run} (tr);
  \draw[ed] (tr) to[bend right=26] node[lbl,above]{clear inputs} (lt);
  \draw[ed] (tr) to[bend left=26]  node[lbl,below]{unload test}  (li);
\end{tikzpicture}
}
\caption{The kernel module's \reg{ioctl} state machine. A test case and inputs load
independently. A run is possible only in \textsc{Ready}. And \textsc{Traced} holds the per-input
measurements, which a re-run reproduces.}
\label{fig:port-execfsm}
\end{figure}

\section{The input format}
\label{sec:port-input}

An input is a self-describing \emph{Revizor Extensible Input File} (\reg{.reif}): a fixed header,
a table of typed section descriptors, and the section payloads (\autoref{fig:port-reif}). The
reader locates each section by its type through the descriptor table rather than by a fixed
offset, so new per-input state can be added without breaking the layout, and an unknown section
type is skipped. The required sections hold the sandbox memory contents and the initial register
values, including the PSTATE condition flags, whose four NZCV bits the fuzzer varies independently.
Optional sections carry further initial state, such as custom PAC keys, an initial MTE tag memory,
and per-page page-table overrides. This extensibility is what later lets precondition
fuzzing attach a \emph{relocation} (code-sealing) section to an input (\autoref{chap:precond})
without changing the file format.

\begin{figure}[t]
\centering
\fitwidth{\begin{tikzpicture}[
  font=\footnotesize,
  blk/.style={draw, rounded corners=2.5pt, align=left, inner sep=4pt, thick, text width=46mm},
  hdr/.style={blk, fill=laneA!12, draw=laneA!60!black},
  dsc/.style={blk, fill=probe!10, draw=probe!55!black},
  pay/.style={blk, fill=laneB!12, draw=laneB!55!black},
  cat/.style={draw, rounded corners=2.5pt, align=left, inner sep=5pt, thick,
              fill=black!2, draw=black!45},
  sz/.style={font=\scriptsize, text=black!55, align=center},
  ar/.style={-{Stealth[length=2mm,width=1.8mm]}, thick, draw=black!65},
]
  % ===== left: file layout =====
  \node[hdr] (h) {%
    {\bfseries header}\hfill{\scriptsize\textcolor{black!55}{48\,B, fixed}}\\[3pt]
    {\ttfamily\scriptsize
     \begin{tabular}{@{}l@{\hspace{2pt}}l@{}}
     magic      & : u64\\
     version    & : u64\\
     header\_len & : u64\\
     n\_sections & : u64\\
     flags      & : u64\\
     total\_len  & : u64\\
     \end{tabular}}};
  \node[dsc, below=2.4mm of h] (d) {%
    {\bfseries section table}\hfill{\scriptsize\textcolor{black!55}{32\,B $\times$ \texttt{n\_sections}}}\\[3pt]
    {\ttfamily\scriptsize
     \begin{tabular}{@{}l@{\hspace{2pt}}l@{}}
     type   & : u64\\
     flags  & : u64\\
     offset & : u64\\
     length & : u64\\
     \end{tabular}}};
  \node[pay, below=2.4mm of d] (p) {%
    {\bfseries section payloads}\\[2pt]
    {\scriptsize\textcolor{black!55}{each begins at its descriptor's \texttt{offset}}}};
  % indirection: a descriptor locates its payload
  \draw[ar] (d.east) to[out=0,in=0,looseness=1.8]
     node[sz, right=7pt, pos=0.5]{\texttt{offset},\\ \texttt{length}} (p.east);

  % ===== right: the section catalog =====
  \node[cat, anchor=north west] (c) at ($(h.north east)+(1.75,0)$) {%
    {\bfseries section types}\quad{\scriptsize\textcolor{probe!55!black}{req}uired / \textcolor{black!50}{opt}ional}\\[3pt]
    \scriptsize
    \begin{tabular}{@{}r@{\ \ }l@{\ \ }l@{\ \ }l@{}}
    \ttfamily 0x01 & \ttfamily MEMORY\_MAIN   & \textcolor{probe!55!black}{req} & sandbox main region\\[1pt]
    \ttfamily 0x02 & \ttfamily MEMORY\_FAULTY & \textcolor{probe!55!black}{req} & sandbox faulty region\\[1pt]
    \ttfamily 0x03 & \ttfamily GPR           & \textcolor{probe!55!black}{req} & \texttt{x0..x5}, \reg{PSTATE} flags, \texttt{sp}\\[1pt]
    \ttfamily 0x04 & \ttfamily SIMD          & \textcolor{black!50}{opt} & \texttt{v0..v7} (reserved)\\[1pt]
    \ttfamily 0x05 & \ttfamily PAC\_KEYS      & \textcolor{black!50}{opt} & per-input \reg{PAC} keys\\[1pt]
    \ttfamily 0x06 & \ttfamily MTE\_TAGS      & \textcolor{black!50}{opt} & one 4-bit tag per 16-byte granule\\[1pt]
    \ttfamily 0x07 & \ttfamily CODE\_RELOC    & \textcolor{black!50}{opt} & word splices into the test body (sealing)\\[1pt]
    \ttfamily 0x09 & \ttfamily PTE\_SETTINGS  & \textcolor{black!50}{opt} & per-page descriptor overrides\\
    \end{tabular}};
  % the descriptor's type names one catalog row
  \draw[ar] (d.north east) to[out=20,in=180] ($(c.west)+(0,-0.15)$);
  \node[sz, anchor=south west] at ($(d.north east)+(0.1,0.02)$) {\texttt{type} selects};
\end{tikzpicture}
}
\caption{The REIF layout. A 48-byte header fixes the magic, version, section count, and total
length, and is followed by one 32-byte descriptor per section, each giving a section's type, byte
offset, and length. The reader finds a section by its type, so the payloads may appear in any
order and an unknown type is skipped. The sandbox memory and the initial registers are required,
and the rest are optional per-input state.}
\label{fig:port-reif}
\end{figure}

\section{Remote and headless fuzzing}
\label{sec:port-remote}

Arm cores commonly ship integrated into phones and embedded devices rather than as parts one
can place in a workstation, so the target is reached over a link. Running the whole fuzzer on
the device is impractical: generating and analyzing large campaigns needs compute, memory, and
a software stack those devices lack. We therefore keep generation, contract modeling, and
analysis on a capable host, and leave on the device hosting the CPU under test only the
execution and measurement that need the silicon. A small on-device helper forwards requests to
the kernel module, so a single host can drive many heterogeneous targets without modification.
Because a network round trip per measurement would dominate the run, a campaign window
aggregates many test cases and their inputs into one streamed super-batch, and the device
returns all traces in a single response.

\begin{figure}[ht]
\centering
\fitwidth{\begin{tikzpicture}[
  font=\footnotesize,
  chip/.style={draw=black!45, rounded corners=2pt, fill=black!4, inner sep=3pt,
               minimum width=34mm, minimum height=5.5mm, align=center, font=\scriptsize},
  mlab/.style={font=\footnotesize\bfseries},
  ar/.style={-{Stealth[length=2.6mm,width=2.2mm]}, very thick, draw=black!72},
  flow/.style={-{Stealth[length=1.8mm,width=1.5mm]}, semithick, draw=black!50},
  lbl/.style={font=\scriptsize, text=black!60, align=center},
  sub/.style={font=\scriptsize, text=black!55},
]
  % ===== host (fuzzing machine) =====
  \node[mlab] (htitle) at (0,1.55) {fuzzing host};
  \node[chip, fill=laneA!14, draw=laneA!55!black] (hg) at (0,0.8)   {generator};
  \node[chip, fill=probe!12, draw=probe!55!black]  (hc) at (0,0.1)  {contract executor (CE)};
  \node[chip, fill=laneB!16, draw=laneB!55!black]  (ha) at (0,-0.6) {analyzer};
  \node[draw=black!55, rounded corners=4pt, thick, fit=(htitle)(hg)(hc)(ha),
        inner sep=4pt] (host) {};
  \node[sub, below=1pt of host] {generation, contract modeling, analysis};

  % ===== target device =====
  \node[mlab] (dtitle) at (7.6,1.55) {target device};
  \node[chip] (dh) at (7.6,0.8)  {on-device helper};
  \node[chip] (dk) at (7.6,0.1)  {kernel module};
  \node[chip, fill=culprit!10, draw=culprit!55!black] (dc) at (7.6,-0.6) {CPU under test};
  \node[draw=black!55, rounded corners=4pt, thick, fit=(dtitle)(dh)(dk)(dc),
        inner sep=4pt] (dev) {};
  \node[sub, below=1pt of dev] {execution and measurement on the silicon};
  \draw[flow] (dh) -- (dk);
  \draw[flow] (dk) -- (dc);

  % ===== link =====
  \draw[ar] ([yshift=4mm]host.east) -- node[lbl, above]{super-batch:\\ test cases $+$ inputs}
        ([yshift=4mm]dev.west);
  \draw[ar] ([yshift=-4mm]dev.west) -- node[lbl, below]{hardware traces\\ (one response)}
        ([yshift=-4mm]host.east);
\end{tikzpicture}
}
\caption{Remote fuzzing. Generation, contract modeling, and analysis stay on the fuzzing host,
while only the kernel module and the CPU under test sit on the target device. The host streams a
super-batch of test cases and their inputs to a small on-device helper, and the device returns
every hardware trace in one response.}
\label{fig:port-remote}
\end{figure}

\medskip
The port reproduces x86 \revizor{}'s methodology on AArch64 and extends it to the Arm security
extensions, but it inherits two limitations of that methodology. Its validation step can
discard genuine leakage that reaches the L1 data cache only indirectly, and its generator
cannot produce the precondition violations behind attacks such as PACMAN and TikTag. The next
two chapters remove these limitations in turn.
